\documentclass[11pt,reqno]{amsart}
\usepackage[localtheorems]{raeez-math-template}

\usepackage{array}
\usepackage{tikz-cd}

\newtheorem{theorem}{Theorem}[section]
\newtheorem{proposition}[theorem]{Proposition}
\newtheorem{lemma}[theorem]{Lemma}
\newtheorem{corollary}[theorem]{Corollary}
\theoremstyle{definition}
\newtheorem{definition}[theorem]{Definition}
\newtheorem{construction}[theorem]{Construction}
\newtheorem{example}[theorem]{Example}
\theoremstyle{remark}
\newtheorem{remark}[theorem]{Remark}

\providecommand{\Ran}{\mathrm{Ran}}
\providecommand{\Sym}{\mathrm{Sym}}
\providecommand{\FactD}{\mathrm{FactD}}
\providecommand{\Alg}{\mathrm{Alg}}
\providecommand{\Coalg}{\mathrm{Coalg}}
\providecommand{\Fin}{\mathrm{Fin}}
\providecommand{\sdim}{\mathrm{sdim}}
\providecommand{\CE}{\mathrm{CE}}
\providecommand{\Bar}{\mathrm{Bar}}
\providecommand{\Cobar}{\mathrm{Cobar}}
\providecommand{\den}{\mathrm{den}}
\providecommand{\id}{\mathrm{id}}
\providecommand{\unit}{\mathbf{1}}
\providecommand{\Hd}{\boxtimes_{\mathrm{H}}}
\providecommand{\Cy}{\boxtimes_{\mathrm{C}}}

\title[Locality and bar denominators]
{The aligned-decomposition locality law and the bar denominator product}

\author{Raeez Lorgat}

\date{July 2026}

\begin{document}

\begin{abstract}
Let $X$ be a smooth complex algebraic curve and let $A$ be an augmented
$E_1$-algebra in factorisation right $D$-modules on $X$. Two structures
on the Ran category, the pointwise coefficient tensor and the chiral
convolution, do not interchange. The canonical comparison between them
is a retraction with an idempotent whose image is the
common-decomposition summand, and strong duoidality holds exactly on the
windows where that idempotent is the identity. At sheaf level the
corresponding exact statement is a criterion: an $E_1$ multiplication in
the ambient Ran category is an algebra in $\FactD(X)$ if and only if
every disjointness-locus square commutes, and those squares occupy the
common-decomposition sector.

For a root-graded Lie superalgebra with finite root spaces, the completed
graded Euler supercharacter of the ordered bar of $U(\mathfrak n_+)$ is
the product side of the Weyl--Kac--Borcherds denominator identity. The
positive Witt algebra gives $\prod_{n\ge1}(1-q^n)$, and the
Gritsenko--Nikulin algebra of $\phi_{0,1}$ gives
$64^{-1}\Delta_5(2Z)$. The Witt specialisation is rigid: Goncharova's
theorem makes the cohomology two-dimensional in every positive degree,
concentrated in the pentagonal weights $(3n^2\pm n)/2$, and Euler's
pentagonal number theorem accounts for the Euler character degree by
degree with no residual freedom. Two circulating dimension formulas for
the bar cohomology of the Virasoro and $\mathfrak{sl}_2$ families are
incompatible with this rigidity; their shared discriminant is an
elementary Motzkin--Riordan identity.
\end{abstract}

\maketitle

\section{The two products on the Ran category}

Let $X$ be a smooth complex algebraic curve, $\Fin^{\mathrm{surj}}$ the
category of nonempty finite sets and surjections, and $\FactD(X)$ the
category of factorisation right $D$-modules: systems $A_I$ on $X^I$ with
multiplicative equivalences on every disjointness locus together with a
filtered extension across the collision diagonals. An
\emph{ordered chiral algebra} is an object
\[
  A \in \Alg^{\mathrm{aug}}_{E_1}\bigl(\FactD(X), \otimes^!\bigr),
\]
an augmented $E_1$-algebra for the pointwise coefficient tensor. The
chiral coefficient records collision on $X$; the internal $E_1$
multiplication records ordered fusion along a transverse oriented line.

The Ran category carries a second monoidal structure, the chiral
convolution assembling disjoint supports. A statement that these two
products interchange, so that an ordered chiral algebra is a bimonoid for
a duoidal structure, asserts an isomorphism where only a retraction
exists. The following two sections isolate the exact content: a
combinatorial idempotent in the additive category of linear species, and
a family of commuting squares at sheaf level.

\section{Aligned decompositions of linear species}

Let $\mathcal V$ be a stable symmetric monoidal category with finite
direct sums. A \emph{linear species} is a functor
$F\colon \Fin^{\simeq}\to\mathcal V$. Species carry two products: the
Hadamard product, componentwise,
\[
  (F \Hd G)(I) = F(I)\otimes G(I),
\]
and the Cauchy product, Day convolution for disjoint union,
\[
  (F \Cy G)(I) = \bigoplus_{I = I_1\sqcup I_2} F(I_1)\otimes G(I_2).
\]
These isolate the finite-set combinatorics beneath the pointwise
coefficient tensor and the chiral convolution respectively. Write an
\emph{ordered decomposition} of $I$ for an ordered pair $d=(I_1,I_2)$
with $I=I_1\sqcup I_2$.

For four species $F_1,F_2,G_1,G_2$ the arity-$I$ component of
$(F_1\Cy F_2)\Hd(G_1\Cy G_2)$ is
\[
  \bigoplus_{d,\,e}
  \bigl[F_1(I_1)\otimes F_2(I_2)\bigr]\otimes
  \bigl[G_1(J_1)\otimes G_2(J_2)\bigr],
  \qquad d=(I_1,I_2),\ e=(J_1,J_2),
\]
a sum over \emph{two} independent ordered decompositions, while the
arity-$I$ component of $(F_1\Hd G_1)\Cy(F_2\Hd G_2)$ is
\[
  \bigoplus_{d}
  \bigl[F_1(I_1)\otimes G_1(I_1)\bigr]\otimes
  \bigl[F_2(I_2)\otimes G_2(I_2)\bigr],
\]
a sum over \emph{one}. The second is the $(d,d)$ sector of the first.

\begin{theorem}[Aligned-decomposition splitting]
\label{thm:aligned-splitting}
There are canonical natural transformations
\[
  \iota\colon (F_1\Hd G_1)\Cy(F_2\Hd G_2)
    \longrightarrow (F_1\Cy F_2)\Hd(G_1\Cy G_2),
\]
\[
  \pi\colon (F_1\Cy F_2)\Hd(G_1\Cy G_2)
    \longrightarrow (F_1\Hd G_1)\Cy(F_2\Hd G_2),
\]
satisfying
\[
  \pi\iota = \id,\qquad e := \iota\pi,\qquad e^2 = e .
\]
The image of $e$ is the common-decomposition summand. Both
transformations are compatible with the associativity and unit
constraints of each product and with the symmetry of the Hadamard
product.
\end{theorem}

\begin{proof}
At a finite set $I$ both source and target are finite direct sums indexed
as displayed. Define $\iota$ to be the inclusion of the summands indexed
by the diagonal $d\mapsto(d,d)$, composed with the symmetry of
$\mathcal V$ reordering the four tensor factors, and $\pi$ to be the
projection onto those summands. Since the index map $d\mapsto(d,d)$ is
injective, $\pi\iota=\id$; consequently $e=\iota\pi$ is the projector
onto the diagonal-indexed direct summand, so $e^2=e$ and
$\operatorname{im} e$ is as stated.

For three Cauchy factors both coherence composites include, respectively
project onto, the triples carrying one common ordered decomposition, so
the associativity constraints are matched. The unit for the Cauchy
product is the species concentrated at $\emptyset$ and the unit for the
Hadamard product is the constant species at the tensor unit; on units the
two index sets coincide and $\iota,\pi$ are identities. Symmetry of the
Hadamard product exchanges the two blocks of the pair $(d,d)$, which the
diagonal index set is stable under. Each identity is checked
componentwise on finite direct sums.
\end{proof}

The theorem is an exact statement in the additive category of species. It
identifies the projector selected by using one common decomposition
simultaneously in both convolution factors, and it shows the failure of
interchange is precisely the presence of the off-diagonal sector
$(d,e)$, $d\ne e$.

\begin{definition}[Aligned species window]
\label{def:aligned-window}
A full replete subcategory $\mathcal S$ of linear species is
\emph{aligned} when it contains both units, is closed under $\Hd$ and
$\Cy$, the idempotent $e$ of
Theorem~\ref{thm:aligned-splitting} acts as the identity on every
iterated mixed product of objects of $\mathcal S$, and the inclusions and
projections remain morphisms of $\mathcal S$.
\end{definition}

\begin{corollary}[Strong duoidality on an aligned window]
\label{cor:aligned-duoidal}
On an aligned window $\iota$ and $\pi$ are mutually inverse interchange
isomorphisms, and $(\Hd,\Cy)$ restricts to a strong duoidal pair.
\end{corollary}

\begin{proof}
The condition $e=\id$ gives $\iota\pi=\id$, and $\pi\iota=\id$ holds by
Theorem~\ref{thm:aligned-splitting}. The coherence identities of that
theorem restrict to the window by fullness and repleteness.
\end{proof}

Fixed decomposition types and finite-support conditions supply elementary
aligned windows. Off such a window the interchange is a retraction and
nothing more, so a bimonoid formulation of locality is unavailable and
must be replaced by the criterion of the next section.

\section{The sheaf-level locality criterion}

Fix an ordered decomposition $d=(I_1,I_2)$ of $I$ and let
$j_d\colon U_d\hookrightarrow X^I$ be its disjointness locus.
Factorisation supplies
\[
  j_d^!A_I \simeq j_d^!\bigl(A_{I_1}\boxtimes A_{I_2}\bigr),
\]
and we write $\Delta_d$ for the inverse factorisation equivalence.

\begin{theorem}[Locality criterion]
\label{thm:locality}
Let $A$ be a factorisation right $D$-module and let
$m\colon A\otimes^! A\to A$ be an $E_1$ multiplication in the ambient Ran
category. Then $m$ defines an $E_1$-algebra in $\FactD(X)$ if and only if
for every ordered decomposition $d$ the square
\[
\begin{tikzcd}[column sep=3.2em, row sep=2.6em]
j_d^!\bigl(A_I\otimes^! A_I\bigr)
  \arrow[r, "j_d^!m"]
  \arrow[d, "\Delta_d\otimes^!\Delta_d"']
& j_d^!A_I \arrow[d, "\Delta_d"] \\
j_d^!\Bigl(\bigl(A_{I_1}\otimes^! A_{I_1}\bigr)\boxtimes
           \bigl(A_{I_2}\otimes^! A_{I_2}\bigr)\Bigr)
  \arrow[r, "m\boxtimes m"]
& j_d^!\bigl(A_{I_1}\boxtimes A_{I_2}\bigr)
\end{tikzcd}
\]
commutes. Under the finite-set indexing functor these squares occupy the
common-decomposition sector of Theorem~\ref{thm:aligned-splitting}.
\end{theorem}

\begin{proof}
A morphism of $\FactD(X)$ is a morphism of the underlying Ran objects
intertwining every factorisation equivalence. Applying that definition to
$m$ produces exactly the displayed family, indexed by ordered
decompositions. Each square uses the same $d$ before and after
multiplication, so under the indexing functor it lies in the sector
indexed by the diagonal $d\mapsto(d,d)$, which is
$\operatorname{im}e$.
\end{proof}

This family is the complete sheaf-level law for ordered locality.
Extension across the partial diagonals remains part of the datum of the
factorisation $D$-module; the indexing projector records only the
separated-support combinatorics.

\begin{proposition}[The ordered bar factorises]
\label{prop:bar-factorises}
Let $B^{\mathrm{ord}}_r(A) = \unit\otimes^! A^{\otimes^! r}\otimes^!\unit$
be the simplicial ordered bar. Every face is built from $m$ or the
augmentation and every degeneracy from the unit, so
Theorem~\ref{thm:locality} applies to each simplicial operator and
$B^{\mathrm{ord}}_\bullet(A)$ is a simplicial object of $\FactD(X)$. Its
realization inherits the factorisation structure and the
deconcatenation coproduct is a morphism of factorisation objects.
\end{proposition}

\begin{proof}
Faces and degeneracies are the structure maps to which the criterion
applies. Realization is computed componentwise. Deconcatenation cuts a
transverse word while a factorisation equivalence separates the chiral
support letter by letter; the two operations commute on every
disjointness locus, which is the morphism condition in $\FactD(X)$.
\end{proof}

Together with the interval identity
$B^{\mathrm{ord}}_X(A)\simeq\unit\otimes^{L}_A\unit$ this places the bar
differential in the transverse direction and keeps the chiral coefficient
uncontracted: no logarithmic residue is asked to produce a bar
differential.

\section{The bar denominator product}

Let
\[
  \mathfrak g = \mathfrak h \oplus
  \bigoplus_{\alpha\in\Gamma\setminus\{0\}} \mathfrak g_\alpha
\]
be a root-graded Lie superalgebra, $\Gamma_+\subset\Gamma$ a positive
monoid, and $\mathfrak n_+ = \bigoplus_{\alpha\in\Gamma_+}\mathfrak
g_\alpha$. Assume every root space is finite dimensional and every degree
admits finitely many decompositions in the chosen $\Gamma_+$-adic
completion. Internal PBW and the bar--Chevalley comparison give
\[
  B^{\mathrm{ord}}\bigl(U(\mathfrak n_+)\bigr)
  \;\simeq\; \CE_*(\mathfrak n_+)
  \;=\; \Sym\bigl(\mathfrak n_+[1]\bigr).
\]

\begin{theorem}[Bar denominator product]
\label{thm:bar-denominator}
The completed graded Euler supercharacter of the ordered bar is
\[
  \chi_\Gamma\Bigl(B^{\mathrm{ord}}\bigl(U(\mathfrak n_+)\bigr)\Bigr)
  \;=\;
  \prod_{\alpha\in\Gamma_+}\bigl(1-e^{-\alpha}\bigr)^{\sdim\mathfrak g_\alpha}.
\]
If $\mathfrak g$ is a Borcherds--Kac--Moody superalgebra with Weyl vector
$\rho$, then
$e^{-\rho}\,\chi_\Gamma\bigl(B^{\mathrm{ord}}(U(\mathfrak n_+))\bigr)$
is the product side of its Weyl--Kac--Borcherds denominator identity.
\end{theorem}

\begin{proof}
Choose a homogeneous basis of each root space. Suspension reverses
parity. An even root vector of root $\alpha$ becomes odd in
$\Sym(\mathfrak n_+[1])$, so it contributes an exterior factor
$1-e^{-\alpha}$ to the Euler supercharacter; an odd root vector becomes
even and contributes a polynomial factor $(1-e^{-\alpha})^{-1}$.
Multiplying over the basis gives the exponent
$\dim\mathfrak g_{\alpha,\bar0}-\dim\mathfrak g_{\alpha,\bar1}
=\sdim\mathfrak g_\alpha$. The finite-decomposition hypothesis makes
every coefficient of the completed product a finite sum. The prefactor
$e^{-\rho}$ is the leading Weyl monomial of the denominator identity.
\end{proof}

The theorem is an object-to-character bridge with the functor named: the
bracket controls the Chevalley differential, and the Euler product
remembers only the signed root spaces. It is not an identification of
the bar object with an automorphic form, and it is not a classification.

\subsection{The Igusa specialisation}

Let $\phi_{0,1}(\tau,z)=\sum_{n\ge0,\ \ell\in\mathbb Z}f(n,\ell)q^ny^\ell$
be the weak Jacobi form of weight $0$ and index $1$ normalized by
$f(0,0)=10$ and $f(0,\pm1)=1$. In the standard cusp chamber its monic
Borcherds product is
\[
  D_5(Z) = q^{1/2}y^{1/2}p^{1/2}
  \prod_{(n,\ell,m)>0}\bigl(1-q^ny^\ell p^m\bigr)^{f(nm,\ell)},
  \qquad \Delta_5(Z) = 64\,D_5(Z),
\]
the second equality being the theta-constant normalization. Gritsenko and
Nikulin \cite{GN97} construct a Borcherds algebra
$\mathfrak g_{\Delta_5}$ with signed root multiplicities
$\sdim(\mathfrak g_{\Delta_5})_{\alpha(n,\ell,m)} = f(nm,\ell)$ and
denominator $\den(\mathfrak g_{\Delta_5}) = D_5(2Z)$.

\begin{corollary}[Igusa bar identity]
\label{cor:igusa}
For the positive nilpotent subalgebra $\mathfrak n_{\Delta_5,+}$,
\[
  e^{-\rho}\,
  \chi_\Gamma\Bigl(B^{\mathrm{ord}}\bigl(U(\mathfrak n_{\Delta_5,+})\bigr)\Bigr)
  \;=\; D_5(2Z) \;=\; 64^{-1}\Delta_5(2Z).
\]
\end{corollary}

\begin{proof}
Insert the Gritsenko--Nikulin root supermultiplicities into
Theorem~\ref{thm:bar-denominator}. The Weyl monomial and the doubled
period variable are those of the type-II denominator chamber, and the
theta-product leading coefficient converts the monic product to
$64^{-1}\Delta_5(2Z)$.
\end{proof}

The root multiplicities are an input, quoted from \cite{GN97}, not a
consequence of Theorem~\ref{thm:bar-denominator}. What the theorem
supplies is the object whose Euler character the product is.

\section{Rigidity of the Witt specialisation}

Let $L_1 = \bigoplus_{n\ge1}\mathbb C\,e_n$ with
$[e_i,e_j]=(j-i)e_{i+j}$, the positive part of the Witt algebra and the
transverse enveloping presentation of the Virasoro family. Every root
space is one dimensional and even, so $\sdim (L_1)_n = 1$ and
Theorem~\ref{thm:bar-denominator} specialises to
\begin{equation}
\label{eq:witt-euler}
  \chi_q\Bigl(B^{\mathrm{ord}}\bigl(U(L_1)\bigr)\Bigr)
  = \prod_{n\ge1}\bigl(1-q^n\bigr).
\end{equation}

\begin{theorem}[Pentagonal rigidity]
\label{thm:pentagonal-rigidity}
The bar cohomology of $U(L_1)$ satisfies
\[
  \dim H^n\Bigl(B^{\mathrm{ord}}\bigl(U(L_1)\bigr)\Bigr) = 2
  \qquad\text{for every } n\ge1,
\]
concentrated in the two weights $(3n^2-n)/2$ and $(3n^2+n)/2$, one
dimension in each. The Euler character \eqref{eq:witt-euler} is
accounted for degree by degree with no residual freedom.
\end{theorem}

\begin{proof}
Under internal PBW the bar of $U(L_1)$ is $\CE_*(L_1)$, whose cohomology
is Goncharova's \cite{Gon73}: $H^n(L_1)$ is two dimensional for $n\ge1$
and lies in the weights $(3n^2\mp n)/2$. Euler's pentagonal number
theorem expands the left side of \eqref{eq:witt-euler} as
$\sum_{m\in\mathbb Z}(-1)^m q^{m(3m-1)/2}$, whose support is exactly
$\{(3n^2\pm n)/2 : n\ge0\}$ with coefficient $(-1)^n$ at each of
$(3n^2-n)/2$ and $(3n^2+n)/2$. Comparing with
$\sum_n(-1)^n\dim H^n_w$ shows every weight of the Euler character is
saturated by a single cohomology class of the degree $n$ predicted by
Goncharova, so no further classes and no cancellation are available.
\end{proof}

Theorem~\ref{thm:pentagonal-rigidity} was verified by exact rational
linear algebra on the weight-graded Chevalley--Eilenberg complex of $L_1$
for $n=1,2,3,4$, with $d^2=0$ asserted on every cochain space rather than
assumed. The computed cohomology occupies the weights
$(1,2)$, $(5,7)$, $(12,15)$, $(22,26)$, one dimension in each, and the
nonzero coefficients of $\prod_{n\ge1}(1-q^n)$ occupy
$0,1,2,5,7,12,15,22,26$ with signs $+,-,-,+,+,-,-,+,+$.

\subsection{Two incompatible dimension formulas}

Two formulas circulate for the bar cohomology of these families:
\[
  \dim H^n = M(n+1)-M(n) = 1,2,5,12,30,76,196,512,1353,3610,\dots
\]
for the Virasoro family, where $M$ denotes the Motzkin numbers, and
\[
  \dim H^n = R(n+3) = 3,6,15,36,91,232,603,1585,\dots
\]
for the $\mathfrak{sl}_2$ family, where $R$ denotes the Riordan numbers;
the coincidence of the discriminant $\sqrt{1-2z-3z^2}$ in the two
generating functions is attributed to Drinfeld--Sokolov reduction.
Neither formula is the bar cohomology of the algebra it names, and the
attribution is unnecessary.

\begin{proposition}[The two sequences are one sequence]
\label{prop:one-sequence}
For all $n\ge0$,
\[
  M(n) = R(n)+R(n+1),
  \qquad
  M(n+1)-M(n) = R(n+2)-R(n).
\]
Hence the Virasoro sequence is $R(n+2)-R(n)$ and the $\mathfrak{sl}_2$
sequence is $R(n+3)$: both are elementary transforms of one sequence, and
their generating functions are rational in the single algebraic function
$\sqrt{1-2z-3z^2}$ for that reason alone.
\end{proposition}

\begin{proof}
Riordan numbers count Motzkin paths with no level step at height zero.
Splitting a Motzkin path of length $n$ according to whether its first
step is level gives $M(n)=R(n)+R(n+1)$, and subtracting consecutive
instances gives the second identity. Both were checked termwise for
$n<24$.
\end{proof}

\begin{proposition}[Neither formula is a cohomology of its algebra]
\label{prop:not-cohomology}
$\dim H^n(\mathfrak{sl}_2;k) = 1,0,0,1$ for $n=0,1,2,3$, and the
Chevalley--Eilenberg chain dimensions are $1,3,3,1$. For an augmented
algebra with $\dim\bar A = d$ the length-$r$ ordered bar term has
dimension $d^r$, so for $\mathfrak{sl}_2$ the bar chain dimensions are
$3,9,27,81,243$. The sequence $3,6,15,36,91$ is neither. For the Virasoro
family, Theorem~\ref{thm:pentagonal-rigidity} gives $\dim H^n = 2$ for
every $n\ge1$, so the sequence $1,2,5,12,30$ is not that cohomology
either.
\end{proposition}

\begin{proof}
The first assertion is Whitehead's lemma for the semisimple Lie algebra
$\mathfrak{sl}_2$, whose cohomology is the invariants
$H^*(\mathfrak{sl}_2;k)=\Lambda(c_3)$ with $c_3$ the Cartan class; it was
recomputed by exact rational elimination on $\Lambda^\bullet
\mathfrak{sl}_2^*$ with $d^2=0$ asserted. The chain counts are the
binomial coefficients $\binom{3}{n}$ and the powers $3^r$. The second
assertion is Theorem~\ref{thm:pentagonal-rigidity}.
\end{proof}

\begin{remark}[Where the Virasoro formula appears to hold]
\label{rem:n2-coincidence}
On degrees $n=1,2,3$ the true dimensions are $2,2,2$ by
Theorem~\ref{thm:pentagonal-rigidity} while the asserted values are
$1,2,5$. The two agree at $n=2$ and nowhere else. A formula tested at
that degree alone reproduces the truth, which is the mechanism by which a
constant sequence can be mistaken for a Motzkin difference.
\end{remark}

\begin{remark}[The type error]
\label{rem:type-error}
Beyond the numerical mismatch the formulas have the wrong type. The bar
cohomology of $U(L_1)$ is bigraded by cohomological degree and weight,
and its total dimension in each degree is finite only because Goncharova
concentrates it in two weights. A formula assigning a single integer to
each degree $n$, with no weight truncation named, does not specify an
invariant of a graded chiral algebra with infinite-dimensional graded
pieces. The sequences $1,2,5,12,30,\dots$ and $3,6,15,36,91,\dots$ are
counts of planar objects; the passage from such a count to a homology
dimension requires a differential, and no quasi-isomorphism to a Motzkin
or Riordan complex has been constructed.
\end{remark}

The honest replacement of both formulas is
Theorem~\ref{thm:bar-denominator}: the invariant of the ordered bar of a
positively graded enveloping algebra that is computable in closed form is
its Euler supercharacter, and that invariant is a denominator product.
For $L_1$ it is $\prod_{n\ge1}(1-q^n)$; the underlying cohomology is two
dimensional in every degree.

\section{The order of dependence}

The constructions above have distinct sources and none determines the
next. The chiral coefficient is the system of $D$-modules with its
factorisation equivalences and diagonal extension. The ordered bar is the
transverse interval amplitude $\unit\otimes^L_A\unit$. An $E_2$
enhancement of the transverse multiplication, or a flat meromorphic
connection on two-dimensional configurations, supplies braid transport. A
cyclic trace supplies the circle object. A clutching-compatible modular
action supplies sewing over stable curves. Scalars occur at specified
functors along
\[
  \FactD(X) \longrightarrow A \longrightarrow B^{\mathrm{ord}}(A)
  \longrightarrow Z^{\mathrm{der}}(A) \longrightarrow \mathrm{Rep}(H)
  \longrightarrow K_0 \longrightarrow \text{character},
\]
and an equality of scalars is explained only by a comparison morphism
between the structured objects preceding them. A central charge, a Hall
polynomial, a determinant, and an automorphic denominator are invariants
at different positions in this sequence; they are not one universal
scalar, and Corollary~\ref{cor:igusa} is an equality of characters
produced by a named functor, not an identification of objects.

\begin{thebibliography}{9}

\bibitem{Bor88}
R.\,E.~Borcherds,
\emph{Generalized Kac--Moody algebras},
J.\ Algebra \textbf{115} (1988), 501--512.

\bibitem{BG96}
A.~Braverman and D.~Gaitsgory,
\emph{Poincar\'e--Birkhoff--Witt theorem for quadratic algebras of
Koszul type},
J.\ Algebra \textbf{181} (1996), 315--328.

\bibitem{FG12}
J.~Francis and D.~Gaitsgory,
\emph{Chiral Koszul duality},
Selecta Math.\ \textbf{18} (2012), 27--87.

\bibitem{Gon73}
L.\,V.~Goncharova,
\emph{Cohomology of Lie algebras of formal vector fields on the line},
Funktsional.\ Anal.\ i Prilozhen.\ \textbf{7} (1973), 6--14.

\bibitem{GN97}
V.~Gritsenko and V.\,V.~Nikulin,
\emph{Siegel automorphic form corrections of some Lorentzian
Kac--Moody Lie algebras},
Amer.\ J.\ Math.\ \textbf{119} (1997), 181--224.

\bibitem{Heu24}
G.~Heuts,
\emph{Koszul duality and a conjecture of Francis--Gaitsgory},
preprint (2024).

\bibitem{Kac90}
V.\,G.~Kac,
\emph{Infinite dimensional Lie algebras},
3rd ed., Cambridge Univ.\ Press, 1990.

\bibitem{Tor25}
T.~Torii,
\emph{On duoidal $\infty$-categories},
J.\ Homotopy Relat.\ Struct.\ \textbf{20} (2025).

\end{thebibliography}

\end{document}
